\section{Существование алгебраического замыкания}

Приведём три доказательства. Первое и третье используют лишь существование простого идеала (принцип ультрафильтра,
\rf{6.4}), второе~--- лемму Цорна и мощности.

\subsection*{Первое доказательство: тензорное произведение полей корня}

Для конечного семейства ненулевых конечномерных коммутативных $F$-алгебр $A_1,\dots,A_n$ тензорное произведение
$A_1\otimes_F\cdots\otimes_FA_n$~--- коммутативная $F$-алгебра с умножением
$(a_1\otimes\cdots\otimes a_n)(b_1\otimes\cdots\otimes b_n)=a_1b_1\otimes\cdots\otimes a_nb_n$. Её размерность равна
произведению размерностей, поэтому она ненулевая. Для конечных $J\subset J'$ имеем
$A_{J'}=A_J\otimes_FA_{J'\setminus J}$, где $A_J=\bigotimes_{j\in J}A_j$, и отображение $A_J\to A_{J'}$, $a\mapsto a\otimes1$,~---
инъективный гомоморфизм алгебр (над полем $V\to V\otimes_FW$, $v\mapsto v\otimes w_0$, инъективно при $w_0\ne0$). Эти
вложения согласованы ($A_J\to A_{J'}\to A_{J''}$ совпадает с $A_J\to A_{J''}$), а система конечных подмножеств направлена
($J,K\subset J\cup K$). Поэтому для произвольного семейства $(A_i)_{i\in I}$ можно склеить все $A_J$, как в \rf{4.25}:
$a\in A_J$ и $b\in A_K$ отождествляются, если их образы в $A_{J\cup K}$ совпадают. Получается ненулевая коммутативная
$F$-алгебра $\bigotimes_FA_i$, порождённая образами канонических гомоморфизмов $A_i\to\bigotimes_FA_i$.

Для каждого унитарного неприводимого $f\in F[X]$ пусть $E_f=F[X]/(f)$~--- его поле корня, и пусть $P$~--- простой
идеал алгебры $A=\bigotimes_FE_f$ (он существует по принципу ультрафильтра; можно взять и максимальный идеал, \rf{6.3}).
Положим $\Om=A/P$. Так как $P\ne A$, кольцо $\Om$ ненулевое, и $F\to\Om$~--- вложение. Это область целостности,
порождённая как $F$-алгебра элементами, алгебраическими над $F$ (образами элементов конечных расширений $E_f$), поэтому
$\Om$~--- поле, алгебраическое над $F$ (\rf{1.32}). Композиция
$E_f\to A\to\Om$~--- гомоморфизм полей, т.\,е.\ вложение; так как у $f$ есть корень в $E_f$, у $f$ есть корень и в
$\Om$. По \rf{6.5} поле $\Om$~--- алгебраическое замыкание $F$.

\subsection*{Второе доказательство: лемма Цорна внутри большого множества}

Если $F$ конечно, оно не более чем счётно, и замыкание построено во введении; поэтому будем считать $F$ бесконечным. Тогда каждое
алгебраическое расширение $F$ имеет мощность $|F|$ (\rf{6.4h}). Возьмём множество $\Xi\supset F$ мощности больше $|F|$
(например, $\Xi=F\cup\mathcal P(F)$: тогда $|\Xi|\ge|\mathcal P(F)|>|F|$ по \rf{6.4d}). Пусть $\mathcal S$~--- множество троек $(E,+_E,\cdot_E)$, где
$F\subset E\subset\Xi$, а $(+_E,\cdot_E)$~--- структура поля на $E$, для которой $F$~--- подполе и $E$ алгебраично над
$F$. Положим $(E,+_E,\cdot_E)\le(E',+_{E'},\cdot_{E'})$, если $E\subset E'$ и $E$~--- подполе $E'$ (операции $E'$,
ограниченные на $E$, совпадают с операциями $E$). Это множество: его элементы~--- подмножества $\Xi$ и операции на них.
Оно непусто (содержит $F$), а у цепи $(E_i)$ есть верхняя грань: объединение $\bigcup E_i$ с операциями, которые для
двух элементов вычисляются в любом $E_i$, содержащем оба; это корректно (операции согласованы), и получается поле в
$\Xi$, алгебраическое над $F$. По лемме Цорна в $\mathcal S$ есть максимальный элемент $\Om$.

Докажем, что $\Om$ алгебраически замкнуто. Пусть, напротив, $g\in\Om[X]$ неприводим и $\deg g>1$. Положим
$K=\Om[X]/(g)$, и пусть $\iota\colon\Om\to K$~--- каноническое вложение. Поле $K$ алгебраично над $F$ (\rf{1.31}), так что
$|K|=|\Om|=|F|$ (\rf{6.4h}). Нам нужно вложение $i\colon K\to\Xi$, тождественное на $\Om$, т.\,е.\ такое, что $i\circ\iota$~--- включение $\Om\subset\Xi$
(коммутативный треугольник ниже); произвольная инъекция $K\to\Xi$ не годится.
\[
\begin{tikzcd}[column sep=large]
\Om\arrow[r,hook,"\iota"]\arrow[dr,hook] & K=\Om[X]/(g)\arrow[d,dashed,cnew,"i"]\\
 & \Xi
\end{tikzcd}
\]
\begin{itemize}[leftmargin=*]
\item \emph{Свободного места хватает:} $|\Xi\setminus\Om|>|F|$. Иначе (\rf{6.4e}) $|\Xi\setminus\Om|\le|F|$ и
$|\Xi|\le|\Om|+|\Xi\setminus\Om|\le|F|+|F|=|F|$ (\rf{6.4g}(a))~--- противоречие с выбором $\Xi$. Поэтому есть инъекция
$j\colon K\setminus\iota(\Om)\to\Xi\setminus\Om$.
\item \emph{Склейка.} Положим $i(\iota(a))=a$ при $a\in\Om$ и $i(x)=j(x)$ при $x\notin\iota(\Om)$. Отображение $i$
инъективно: на каждом из двух кусков оно инъективно, а образы кусков лежат в непересекающихся множествах $\Om$ и
$\Xi\setminus\Om$.
\item \emph{Перенос структуры.} На $S=i(K)\subset\Xi$ определим операции
$s\oplus t=i\bigl(i^{-1}(s)+i^{-1}(t)\bigr)$, $s\odot t=i\bigl(i^{-1}(s)\cdot i^{-1}(t)\bigr)$. Тогда $S$~--- поле, а
$i\colon K\to S$~--- изоморфизм. Для $a,b\in\Om$ имеем $a\oplus b=i(\iota(a)+\iota(b))=i(\iota(a+b))=a+b$ и так же для
умножения, т.\,е.\ $\Om$~--- подполе $S$. Поле $S$ алгебраично над $F$ (изоморфно $K$ над $\Om\supset F$) и отлично от
$\Om$, так как $K\ne\iota(\Om)$.
\end{itemize}
Итак, $\Om<S$ в $\mathcal S$, что противоречит максимальности. Значит, $\Om$ алгебраически замкнуто и, будучи
алгебраическим над $F$, является алгебраическим замыканием $F$.

\begin{center}
\begin{tikzpicture}[>=stealth]
\begin{scope}
\clip (0,0) ellipse (1.7 and 1.05);
\fill[cold!14] (-1.7,-1.1) rectangle (0.15,1.1);
\fill[cnew!14] (0.15,-1.1) rectangle (1.7,1.1);
\end{scope}
\draw (0,0) ellipse (1.7 and 1.05);
\draw[gray] (0.15,-1.05) -- (0.15,1.05);
\node[cold] at (-0.75,0) {$\iota(\Om)$};
\node[cnew,font=\footnotesize,align=center] at (0.9,0) {$K\setminus$\\ $\iota(\Om)$};
\node at (0,-1.35) {$K=\Om[X]/(g)$};
\draw[rounded corners=8pt] (3.6,-1.6) rectangle (9.4,1.6);
\node at (9.1,1.3) {$\Xi$};
\draw[fill=cold!14] (5.0,0.1) ellipse (1.0 and 0.75);
\node[cold] at (5.0,0.1) {$\Om$};
\draw[fill=cnew!14] (7.7,-0.35) ellipse (0.95 and 0.55);
\node[cnew,font=\footnotesize] at (7.7,-0.35) {$j\bigl(K\setminus\iota(\Om)\bigr)$};
\node[font=\footnotesize] at (7.0,1.05) {$\Xi\setminus\Om$};
\draw[->,cold] (-0.6,0.7) to[bend left=25] node[above,font=\footnotesize]{$\iota(a)\mapsto a$} (4.5,0.75);
\draw[->,cnew] (1.0,-0.65) to[bend right=18] node[pos=0.3,below,font=\footnotesize]{$j$} (6.95,-0.7);
\end{tikzpicture}
\end{center}
\figcap{Вложение $i\colon K\to\Xi$ склеено из тождественного отображения на $\Om$ и инъекции $j$ <<новой части>> $K$ в
свободную часть $\Xi\setminus\Om$; её хватает, потому что $|K|=|F|<|\Xi\setminus\Om|$. Затем на $i(K)$ переносится
структура поля.}

Работать приходится внутри фиксированного множества $\Xi$, так как совокупность всех алгебраических расширений $F$ не
является множеством; условие $|\Xi|>|F|$ нужно ради шага со <<свободным местом>>: при малом $\Xi$ (например, $\Xi=F$)
максимальный элемент не обязан быть алгебраически замкнут. Тот же приём (тождественно на старой
части, инъекция новой части в свободное место, перенос структуры) стандартен всюду, где объект строят внутри
фиксированного множества, чтобы применить лемму Цорна.

\subsection*{Третье доказательство: кольцо многочленов от многих переменных}

Рассмотрим кольцо многочленов $R=F[\dots,X_f,\dots]$ от переменных $X_f$, занумерованных унитарными неприводимыми
многочленами $f\in F[X]$ (каждый элемент $R$ зависит от конечного числа переменных), и идеал $I\subset R$, порождённый
многочленами $f(X_f)$.

\emph{$I\ne R$.} Если $1\in I$, то
\[
g_1f_1(X_{f_1})+\cdots+g_nf_n(X_{f_n})=1\qquad\text{в }R
\]
для некоторых $g_k\in R$ и попарно различных унитарных неприводимых $f_k$. Пусть $E$~--- конечное расширение $F$, в котором у каждого
$f_k$ есть корень $\alpha_k$ (например, поле разложения $f_1\cdots f_n$, \rf{2.7}). Гомоморфизм $F$-алгебр $R\to E$,
$X_{f_k}\mapsto\alpha_k$, $X_f\mapsto0$ при $f\notin\{f_1,\dots,f_n\}$, переводит это равенство в $0=1$~---
противоречие.

Пусть $P\supset I$~--- простой идеал $R$ (существует по принципу ультрафильтра, применённому к $R/I\ne0$; можно взять и
максимальный идеал, \rf{6.3}), и $\Om=R/P$. Так как $P\ne R$, кольцо $\Om$ ненулевое и $F\to\Om$~--- вложение. Это
область целостности, порождённая над $F$ образами $\bar X_f$, которые алгебраичны над $F$ (так как $f(\bar X_f)=0$);
по \rf{1.32} $\Om$~--- поле, алгебраическое над $F$.
Каждый унитарный неприводимый $f$ имеет в $\Om$ корень $\bar X_f$, значит, каждый непостоянный многочлен из $F[X]$ имеет
корень в $\Om$, и по \rf{6.5} $\Om$~--- алгебраическое замыкание $F$.
